\documentclass[11pt]{amsart}
\usepackage{amsmath,amssymb,amsthm}
\usepackage[margin=1in]{geometry}
\usepackage[colorlinks=true,linkcolor=blue,citecolor=blue,urlcolor=blue]{hyperref}
\newtheorem{theorem}{Theorem}[section]
\newtheorem{lemma}[theorem]{Lemma}
\newtheorem{proposition}[theorem]{Proposition}
\newtheorem{corollary}[theorem]{Corollary}
\theoremstyle{remark}\newtheorem{remark}[theorem]{Remark}
\newcommand{\E}{\mathbb{E}}
\title[A prophet secretary instance below 0.7235]{A prophet secretary instance with ratio below 0.7230}
\author{Yitzchak Shmalo}
\begin{document}
\begin{abstract}
In the prophet secretary problem, independent nonnegative random variables with known distributions arrive in
uniformly random order, and a gambler who sees each value as it arrives may stop once and keep it. The best constant
$\alpha$ such that the gambler can always obtain $\alpha$ times the expected maximum lies between $0.688$ and $0.7235$,
the upper bound being due to Giambartolomei, Mallmann-Trenn and Saona. We give an explicit instance with $461$
variables on which the optimal gambler obtains less than $0.722934$ times the expected maximum, so that
$\alpha<0.72294$. The instance consists of $400$ rare prizes, two groups of $30$ medium items with two different
values, and one deterministic item. Its optimal value is computed exactly by a dynamic programme over counts of
remaining items, evaluated in interval arithmetic. In the limit of many prizes the same family is described by an
ordinary differential equation, whose numerical solution suggests that two medium values give about $0.72274$ and
three about $0.72268$.
\end{abstract}
\maketitle

\section{Introduction}

Let $X_1,\dots,X_n$ be independent nonnegative random variables with known distributions. In the prophet secretary
problem the variables are presented in a uniformly random order $\sigma$; at step $t$ the gambler learns the index
$\sigma(t)$ and the value $X_{\sigma(t)}$, and either stops and receives $X_{\sigma(t)}$ or rejects it for good. Let
$V(X_1,\dots,X_n)$ be the largest expected reward of a stopping rule adapted to this information, and let
\[
\alpha=\inf\frac{V(X_1,\dots,X_n)}{\E\max_i X_i},
\]
the infimum over all $n$ and all instances with $0<\E\max_i X_i<\infty$.

Esfandiari, Hajiaghayi, Liaghat and Monemizadeh \cite{ehlm} showed $\alpha\ge 1-1/e$, and the best lower bound is now
$0.688$, due to Chen, Huang, Li and Tang \cite{chlt}. On the hardness side, Correa, Saona and Ziliotto \cite{csz} showed
$\alpha\le\sqrt3-1\approx0.732$, Bubna and Chiplunkar \cite{bc} obtained $0.7254$ by a computer-assisted argument (their
bound holds even for a gambler who knows the order in advance), and Giambartolomei, Mallmann-Trenn and Saona \cite{gms}
proved $\alpha<0.7235$ with an instance whose ratio they computed in closed form in a limit. We improve the upper bound.

\begin{theorem}\label{thm:main}
There is an instance with $461$ variables, each taking one positive value or $0$, whose optimal prophet secretary value
satisfies
\[
0.7229334600126004776\le\frac{V}{\E\max_i X_i}\le0.7229334600126004777 .
\]
In particular $\alpha<0.72294$.
\end{theorem}

The instance, called $I(400,30)$ below, has the shape of the instance of \cite{gms}. There are many independent prizes,
each so rare that with high probability at most one of them appears, so that their total mass behaves like a stream of
reward that arrives at a constant rate while the gambler waits. There is one item that is certainly present and is worth
less than the prize mass, and a cloud of medium items, each present with small probability, worth slightly more than the
deterministic item. The gambler who meets the deterministic item early must decide whether to give up the remaining
prize stream; the medium items make both waiting and stopping costly. In \cite{gms} the medium items all have one value.
Taking two values, as anticipated in the conclusion of \cite{gms}, lowers the ratio.

The value of the instance is computed by the standard recursion for the optimal stopping rule under random order
(Section~\ref{sec:dp}). Because items in a group are identically distributed, the recursion depends only on how many items
of each group are left, which reduces the $2^{461}$ subsets to $770{,}722$ states (Section~\ref{sec:cert}). The
recursion uses only arithmetic operations and maxima, so it can be evaluated in interval
arithmetic with outward rounding, and the result is an enclosure of the exact ratio. Section~\ref{sec:limit} describes
the limit of the family as the number of prizes grows; the values reported there are numerical and are not part of the
theorem.

\section{The recursion}\label{sec:dp}

For a set $S$ of indices let $V(S)$ be the optimal expected reward of the game played on the variables $(X_i)_{i\in S}$
alone, in uniformly random order, with $V(\emptyset)=0$.

\begin{lemma}\label{lem:dp}
For nonempty $S$,
\[
V(S)=\frac1{|S|}\sum_{i\in S}\E\max\bigl(X_i,V(S\setminus\{i\})\bigr).
\]
\end{lemma}

\begin{proof}
The first index presented is uniform on $S$. Given that it is $i$, the remaining order is uniform on the orderings of
$S\setminus\{i\}$, and the remaining variables are independent of $X_i$ and of the order. Hence after observing $X_i$
the gambler who continues faces exactly the game on $S\setminus\{i\}$, whose optimal value is $V(S\setminus\{i\})$
whatever has been observed, and the optimal decision is to stop if and only if $X_i>V(S\setminus\{i\})$. Averaging over
$i$ gives the formula; a backward induction over $|S|$ shows that the rule defined in this way is optimal among all
stopping rules adapted to the observed indices and values.
\end{proof}

\begin{lemma}\label{lem:counts}
Suppose the index set is partitioned into groups $G_1,\dots,G_r$, the variables in each group being identically
distributed, say $X_i\overset{d}{=}Y_g$ for $i\in G_g$. Then $V(S)$ depends on $S$ only through the counts
$c_g=|S\cap G_g|$, and writing $V(c)$ for it,
\[
V(c)=\frac1{|c|}\sum_{g:\,c_g>0}c_g\,\E\max\bigl(Y_g,V(c-e_g)\bigr),\qquad |c|=c_1+\dots+c_r,\quad V(0)=0 .
\]
\end{lemma}

\begin{proof}
Induction on $|S|$ using Lemma~\ref{lem:dp}: the terms with $i\in S\cap G_g$ are all equal, since $X_i$ has the law of
$Y_g$ and $S\setminus\{i\}$ has counts $c-e_g$.
\end{proof}

If $Y_g$ equals $v_g>0$ with probability $\pi_g$ and $0$ otherwise, then
$\E\max(Y_g,w)=\pi_g\max(v_g,w)+(1-\pi_g)w$ for $w\ge0$. If independent variables $X_i$ equal $v_i$ with
probability $\pi_i$ and $0$ otherwise, and $v_1\ge v_2\ge\cdots$, then
\begin{equation}\label{eq:emax}
\E\max_i X_i=\sum_i v_i\pi_i\prod_{j<i}(1-\pi_j).
\end{equation}

\section{The instance and the certificate}\label{sec:cert}

Fix the parameters
\[
a=0.823,\quad b_1=1.28233,\quad p_1=0.2065,\quad b_2=1.38487,\quad p_2=0.17698,\quad q=10^{-9},
\]
and integers $k,m\ge1$. The instance $I(k,m)$ has four groups:
\begin{itemize}
\item $k$ prizes, each equal to $1$ with probability $q$ and to $0$ otherwise;
\item $m$ items, each equal to $B_1=b_1kq$ with probability $p_1/m$ and to $0$ otherwise;
\item $m$ items, each equal to $B_2=b_2kq$ with probability $p_2/m$ and to $0$ otherwise;
\item one item equal to $A=akq$.
\end{itemize}
All values of the medium and deterministic items are measured in units of the total prize mass $kq$; for $k=400$ these
values are of order $4\cdot10^{-7}$, far below the prize value $1$. Since $1>B_2>B_1>A$, \eqref{eq:emax} gives
\[
\E\max=1-(1-q)^k+(1-q)^k\Bigl[B_2\bigl(1-(1-\tfrac{p_2}{m})^m\bigr)+(1-\tfrac{p_2}{m})^m\Bigl(B_1\bigl(1-(1-\tfrac{p_1}{m})^m\bigr)
+(1-\tfrac{p_1}{m})^mA\Bigr)\Bigr].
\]
By Lemma~\ref{lem:counts} the optimal value is $V(k,m,m,1)$, computed over the states $(j,l_1,l_2,s)$ with
$0\le j\le k$, $0\le l_1,l_2\le m$, $s\in\{0,1\}$, of which there are $2(k+1)(m+1)^2$; for $k=400$, $m=30$ this is
$770{,}722$.

\begin{proof}[Proof of Theorem~\ref{thm:main}]
Take the instance $I(400,30)$, which has $400+30+30+1=461$ variables. The recursion of Lemma~\ref{lem:counts} was
evaluated over all $770{,}722$ states in interval arithmetic (the \texttt{iv} context of mpmath with a $100$-bit
mantissa), with each decimal parameter replaced by an interval containing it. Every operation in the recursion is an
arithmetic operation or a maximum of two numbers; interval arithmetic with outward rounding returns an interval
containing the exact result of an arithmetic operation applied to any points of its argument intervals, and
$[x_1,y_1]\vee[x_2,y_2]=[\max(x_1,x_2),\max(y_1,y_2)]$ contains $\max(x,y)$ for every $x\in[x_1,y_1]$, $y\in[x_2,y_2]$.
By induction over the states, each computed interval contains the exact value of its state. The same holds for the
closed form of $\E\max$. The run returned
\[
\frac{V}{\E\max}\in[\,0.722933460012600477658647124,\;0.722933460012600477658647127\,],
\]
which contains the stated interval.
\end{proof}

\begin{remark}
The computation takes about two and a half minutes on one core. Three further checks were made, none of which the proof
uses. The recursion over counts agrees with the recursion of Lemma~\ref{lem:dp} over all $2^n$ subsets, run separately
with every item distinct, for $I(6,2)$ ($n=11$, equal to the last bit in double precision) and $I(8,3)$ ($n=15$,
difference $1.1\cdot10^{-16}$). Double precision gives $0.7229334600126$ for $I(400,30)$, in agreement with the
certificate. Replacing $q=10^{-9}$ by $q=10^{-12}$ changes the ratio of $I(400,30)$ by $7\cdot10^{-8}$, so the tiny prize
probability enters only through the prize mass $kq$, as the construction intends.
\end{remark}

At the same parameters the ratio decreases as the instance grows:
\[
\begin{array}{c|cccc}
(k,m) & (100,10) & (200,20) & (400,30) & (800,40)\\\hline
V/\E\max & 0.7234029 & 0.7230763 & 0.7229335 & 0.7228621
\end{array}
\]
(double precision). Every row is a valid instance; the theorem uses the third because it is the smallest one we
certified. The successive differences roughly halve as $k$ doubles.

\section{The many-prize limit}\label{sec:limit}

The following description of the limit is heuristic, and the numbers in this section are numerical solutions of an
ordinary differential equation; they are not certified and are not claimed as bounds.

Let $k\to\infty$ with the prize mass $kq\to0$, measure rewards in units of $kq$, and let the number of medium items of
each value grow. Rescale time so that the game runs over $[0,1]$. Then the prizes become a stream that pays reward at
rate $1$ for as long as the gambler waits (a single prize is always accepted, and to first order at most one appears),
the medium items of value $b_j$ arrive as a Poisson process of rate $p_j$, and the deterministic item arrives at a
uniform time. Writing $u$ for the time remaining, the value $W(u)$ of the game after the deterministic item has passed and
the value $\Phi(u)$ before it solve
\[
W'=1+\sum_jp_j(b_j-W)^+,\quad W(0)=0,\qquad
\Phi'=1+\sum_jp_j(b_j-\Phi)^++\frac{\max(a,W)-\Phi}{u},\quad \Phi(0)=a,
\]
and the limiting ratio is $\Phi(1)$ divided by $1+\E\max(a,Z)$, where $Z$ is the largest medium value that appears. With
no medium items, $W(u)=u$ and $\Phi(1)=1+a^2/2$ for $a\le1$, and the ratio $(1+a^2/2)/(1+a)$ is minimised at $a=\sqrt3-1$
with value $\sqrt3-1$, the instance of \cite{csz}. With one medium value this is the family of \cite{gms}.

Solving these equations numerically (an adaptive Runge--Kutta method of order $8$, relative tolerance $10^{-12}$, at
parameters found by Nelder--Mead minimisation of the limit ratio) gives the following values.
\begin{center}\small
\begin{tabular}{l|l|l|l}
medium values & $a$ & $(b_j,p_j)$ & limit ratio\\\hline
none & $\sqrt3-1$ & & $0.73205081$\\
one, parameters of \cite{gms} & $0.789$ & $(1.24,0.421)$ & $0.72348603$\\
one, minimised & $0.8204$ & $(1.3304,0.3717)$ & $0.72302401$\\
two, parameters of $I(k,m)$ & $0.823$ & $(1.28233,0.2065)$, $(1.38487,0.17698)$ & $0.72275062$\\
two, minimised & $0.82160$ & $(1.27865,0.20840)$, $(1.38431,0.18833)$ & $0.72274261$\\
three, minimised & $0.82178$ & $(1.26006,0.14352)$, $(1.33157,0.13374)$, & $0.72268397$\\
 & & $(1.40321,0.12514)$ &
\end{tabular}
\end{center}
The second row agrees with the bound $0.7235$ of \cite{gms}. The values of $I(k,m)$ in Section~\ref{sec:cert} approach the
fourth row from above, the remaining gap at $k=800$ being due to the finite number of medium items. With three values the
optimal medium values are nearly evenly spaced, with nearly equal rates, which suggests that a continuous band of medium
values is the extremal configuration within this family. The gains from adding a value decrease quickly
($2.8\cdot10^{-4}$ from one to two, $5.9\cdot10^{-5}$ from two to three).

\section*{Acknowledgments and funding}

This research was supported by the European Research Council (ERC) under the European Union's Horizon Europe research and innovation programme (grant agreement No.~101041711), by the Simons Foundation as part of the Collaboration on the Mathematical and Scientific Foundations of Deep Learning, by Heights Labs, by Convex Nexus Capital, by the Israel Science Foundation (grant number 2258/19), and by the Israel Science Foundation (ISF Grant 4101/25).

\section*{Declaration of competing interest}

The author holds equity in Heights Labs and in Convex Nexus Capital and has a management role at Convex Nexus Capital. These entities are included in the funding declaration above.

\section*{Declaration of generative AI and AI-assisted technologies}

Anthropic Claude AI systems were used to develop the mathematical derivations, write and run the numerical checks, perform internal consistency and literature reviews, and write and prepare this manuscript and its publication files. The author selected the problem and directed the work; the author did not carry out the mathematics or the writing. Responsibility for submission or publication rests with the author. The AI-assisted internal reviews, compilation checks and rendered-page inspections are not external peer review or formal verification.

\section*{Data and code availability}

The certificate of Theorem~\ref{thm:main}, the checks of Section~\ref{sec:cert} and the numerical solutions of
Section~\ref{sec:limit} are produced by the scripts in the \texttt{code} directory of the repository
\[
\texttt{github.com/yspennstate/prophet-secretary-hardness}.
\]

\begin{thebibliography}{10}
\bibitem{bc} A.~Bubna and A.~Chiplunkar, \emph{Prophet inequality: order selection beats random order}, Proc. 24th ACM
Conf. Economics and Computation (EC 2023), arXiv:2211.04145.
\bibitem{chlt} Z.~Chen, Z.~Huang, D.~Li and Z.~G.~Tang, \emph{Prophet secretary and matching: the significance of the
largest item}, Proc. 2025 ACM-SIAM Symp. Discrete Algorithms (SODA 2025), arXiv:2411.01191.
\bibitem{csz} J.~Correa, R.~Saona and B.~Ziliotto, \emph{Prophet secretary through blind strategies}, Math. Program.
\textbf{190} (2021), 483--521, arXiv:1807.07483.
\bibitem{ehlm} H.~Esfandiari, M.~Hajiaghayi, V.~Liaghat and M.~Monemizadeh, \emph{Prophet secretary}, SIAM J. Discrete
Math. \textbf{31} (2017), 1685--1701.
\bibitem{gms} G.~Giambartolomei, F.~Mallmann-Trenn and R.~Saona, \emph{Prophet inequalities: separating random order
from order selection}, arXiv:2304.04024.
\end{thebibliography}

\end{document}
